In what follows, we use the term agent(s)---common in computational social choice literature---interchangeably with point(s)---as used in the clustering literature. 
Let $\Mcal$ denote the set of all metric spaces on finite points.
Let $(U, d) \in \Mcal$ be a metric space with distance function $d: U \times U \rightarrow \RR_{\ge 0}$.
For a subset $S \subseteq U$ of points, we use $d(u,\SolSet)$ to denote $\min_{s \in \SolSet} d(u, s)$. 
We consider the setting where ordinal rankings of all points of $U$ are available. 
Specifically, for each point $v \in U$, a linear order $\succ_v: [n] \rightarrow U$ is known. We write $u \succ_{v} u'$ to indicate that $u$ is closer $v$ than $u'$ in the ranking $\succ_v$. 
Furthermore, we assume that the linear orders are \emph{consistent} with $d$,\ie, for every $v \in U$ and for all $u, u' \in U$, $u \succ_v u'$ implies $d(u, v) \leq d(u', v)$. 
An \emph{ordinal profile} of $U$ is a collection of linear orders of all the points of $U$, which we denote by $\succ_U$.
We say an ordinal profile $\succ_U$ of $U$ is \emph{consistent} with $d$ if all the linear orders in $\succ_U$ are consistent with $d$.
Let $\profile(d)$ denote the set of all ordinal profiles of $U$ that are consistent with $d$.
We first introduce the ordinal fair $k$-center problem and make it precise how social cost is defined in this context.

\begin{definition}[The ordinal fair $k$-center problem (\ordfairkcen)]
\label{def:ordinal-fair-k-center-with-limited-distance-queries}
An instance of the ordinal fair $k$-center problem is defined on a set $U$ of $n$ points from a metric space $(U,d)$ with unknown $d$, an integer $k \geq  1$, a collection $\succ_U=\{\succ_u\}_{u \in U}$ of linear orders that is  consistent with $d$, a collection $\GroupSet= \{\Group_1, \dots, \Group_t : \Group_i \subseteq U\}$ of $t$ subsets of data points 
that form a partition of $U$,
and a vectors of requirements $\alphavec = \{\alpha_1,\dots,\alpha_t\}$, where $\alpha_i \ge 0$ corresponds to the requirement of group $\Group_i$. A set $\SolSet \subseteq U$ of centers is a feasible solution if $|\SolSet| = k$ and $\alpha_i \leq |\SolSet \cap \Group_i|$ for all $i \in [t]$. 
The \emph{social cost} of a solution $\SolSet$ is the maximum distance of any point to $S$, \ie, $\cost(S) = \max_{u \in U} d(u,\SolSet)$. The goal of the ordinal fair $k$-center problem is to find a feasible solution with minimum social cost.
\end{definition}


An instance of the ordinal fair $k$-center problem is denoted as $\Ical = \ordfairkcenins$.
The ordinal $k$-center problem is defined when no fairness constraints are enforced.
The fair $k$-center problem is defined analogously, with the key distinction that the distance function $d$ of the underlying metric space is fully known.  
For brevity, we denote the ordinal $k$-center problem with and without fairness constraints by \ordfairkcen\ and \ordkcen, respectively.

We evaluate the quality of the solution via distortion~\cite{procaccia2006distortion}, defined as the approximation ratio between the social cost of the algorithm's solution 
and the social cost of the optimal solution (computed with access to all pairwise distances).
We first formally define distortion for \ordfairkcen, adopting the notion from Burkhardt et~al.~\cite{burkhardt2024low}.




\begin{definition}[Distortion of Algorithm $\alg$]
    For a metric space $(U,d) \in \Mcal$ and an ordinal profile $\succ_U \in \profile(d)$, let $\Ical_{\succ_U}$ be the collection of instances of the ordinal (fair) $k$-center problem defined on points $U$ and ordinal profile $\succ_U$. 
    Let $\alg(I \mid d)$ denote the solution returned by algorithm $\alg$ on instance $I \in \Ical_{\succ_U}$ with underlying metric $d$. 
    Let $\OptSolSet_{I, d}$ be an optimal solution to $I$ when the underlying metric is $d$.
    The distortion of $\alg$ is defined as:
    \[
    \distortion(\alg)\coloneqq \sup_{\substack{(U,d) \in \Mcal \\ \succ_U \in \profile(d)}} \sup_{I \in \Ical_{d,\succ_U}}  \frac{\cost(\alg(I \mid d))}{\cost(\OptSolSet_{I,d})}
    \]
\end{definition}


Since without access to cardinal information, it is not possible to obtain an algorithm with bounded distortion for \ordkcen~\cite{burkhardt2024low}, we allow the algorithm to make \emph{query access} to $d$: given a pair $u,v\in U$, the algorithm can query the distance $d(u,v)$. 
Hence, our goal is to design an algorithm $\alg$ that uses a limited number of queries while ensuring that the solution quality is provably close to the optimal solution. 


\xhdr{Reduction} 
For simplicity of exposition, we transform an instance $\Ical = (U, k, \succ_U, \GroupSet, \alphavec)$ of the ordinal fair $k$-center problem, where $\GroupSet = \{\Group_1, \dots, \Group_t\}$ and $\alphavec = \{\alpha_1, \dots, \alpha_t\}$, into an equivalent instance $\Ical' = (U', k, \succ_{U'}, \GroupSet', \alphavec')$ with exactly $k$ groups.
Formally, we construct the transformed instance $\Ical'$ as follows. 
Let $r = \sum_{i \in [t]} \alpha_i$. 
For each group $G_i \in \GroupSet$, we make $\alpha_i$ disjoint copies of $G_i$ in $\GroupSet'$ by duplicating each element of $G_i$ precisely $\alpha_i$ times. Next, if $r < k$, then we create $(k-r)$ new groups in $\GroupSet'$ each containing a distinct copy of $U$. We set $\alphavec' = \{1, 1, \dots, 1\}$, so that the fairness constraint requires selecting exactly one center from each group $\Group'_j$ for $j \in [k]$. 
We define the linear orderings in $\succ_{U'}$ as follows. 
For each point $v' \in U'$, let $v \in U$ be the original point of $v'$ (where $v'$ may be $v$ itself or a duplicate copy of $v$). 
We define the linear order $\succ_{v'}$ on $U'$ as an order extension of $\succ_{v}$ on $U$: for any points $a', b' \in U'$ and their corresponding original points $a, b \in U$, we have $a' \succ_{v'} b'$ if and only if $a \succ_{v} b$.
When $a = b$, the ordering of $a'$ and $b'$ is arbitrary.







This transformation is standard in the fair clustering literature and has been used in several prior works~\cite{thejaswi2022clustering, thejaswi2024diversity, gadekar2025fair}, in this transformed instance $\Ical'$, three key properties are relevant for our analysis to hold: ($i$) the cost of optimal solution in $\Ical'$ remains identical to that of the original instance $\Ical$, since the additional points are duplicates; {($ii$) no extra distance queries are required, as duplicate points share identical distances;} and ($iii$) if the algorithm's output for $\Ical'$ contains multiple copies of the same center, we retain a single copy and supplement it with arbitrary 
centers
from the corresponding group to satisfy the fairness constraints, without affecting the theoretical guarantees on the distortion factor.

From now on, we present our algorithmic results for the instance $\Ical'$, and we directly write the instance as $\Ical$ when there is no ambiguity.




\xhdr{Preliminaries}
In what follows, we define the terminology and definitions necessary to present our algorithmic results.
Given a  set $U$ and two subsets $A,B \subseteq U$ of points, we say that $A$ hits $B$ if $A \cap B \neq \emptyset$. More generally, for an integer $i \geq 0$, $A$ hits $B$ $i$ times if $|A \cap B|=i$. For an ordered set $T = (t_1,\dots,t_k) \subseteq U$, and $\ell \in [k]$, we denote by $T_\ell$ the $\ell$-length prefix of $T$, \ie, $T_\ell= (t_1,\dots, t_\ell)$.

To characterize the quality of Gonzalez-type greedy algorithms~\cite{gonzalez1985clustering, burkhardt2024low}, we apply the notion of progressive cover and critical index.
Let $\mathbf{\Pi}^*$ denote the partition of $U$ induced by the optimal solution $S^*$; \ie, each cluster contains all points whose nearest center in $S^*$ is the same.
Note that the $k$ centers $T$ obtained by Gonzalez's classic $k$-center algorithm~\cite{gonzalez1985clustering} form a progressive $2$-cover. 
In this case, the critical index for $T$ is the largest $\ell \in [k]$ such that $T_{\ell}$ hits each part of the partition $\mathbf{\Pi}^*$ at most once, \ie, the largest $\ell$ that satisfies (i) also satisfies (ii) in the later definition. 
This property that the largest $\ell$ that satisfies (i) also satisfies (ii) also holds for the solution obtained by the algorithm of Burkhardt et~al.~\cite{burkhardt2024low} for the ordinal $k$-center problem.

\begin{definition}[Progressive cover and critical index]
Let $\Jcal$ be an instance of \kcenter (or \fairkcenter) and let $S^*$ be an optimal solution with corresponding partition $\mathbf{\Pi}^*$ of $U$. Fix some $\gamma >0$, and consider an ordered set $T=(t_1,\dots,t_k) \subseteq U$ of size $k$. We say $T$ is a \emph{progressive $\gamma$-cover} for $\Jcal$ \wrt $S^*$, if there exists $\ell \in [k]$ such that ($i$) $T_{\ell}$ hits each part of $\mathbf{\Pi}^*$ at most once, and ($ii$) $\cost(T_{\ell}) \leq \gamma\cdot \cost(S^*)$. Furthermore, we say $\ell \in [k]$ as the \emph{critical index} of a progressive $\gamma$-cover $T$ for $\Jcal$  \wrt $S^*$ if $\ell$ is the maximum index such that ($i$) and ($ii$) hold.
\end{definition}

When $S^*$ is implicit from context, we omit saying \wrt $S^*$.
A crucial tool that is used to obtain a feasible (fair) solution from a progressive $\gamma$-cover $T$~\cite{gadekar2025fair,jones2020fair} is to project $T$ onto a fair solution using a matching on a carefully constructed bipartite graph. 
In our paper, we apply the notion of left-perfect matching to define the projection graph. We recall the definition of left-perfect matching in Definition~\ref{definition:left-perfect-matching} in the appendix.
We formally define the projection graph below.
\begin{definition}[Projection graph]
Consider an instance $\Jcal=\fairkcenins$ of \fairkcenter, and let $S \subseteq U, |S|\le k$. Furthermore, let $\ell \in [|S|]$, and $\lambda \in \mathbb{R}_+$. 
We define the \emph{$(\ell, \lambda)$-projection graph}  $H^{\ell}_{\lambda} = (\InitSolSet_{\ell} \cup \MultiGroupSet, E_{\lambda})$, where there is an edge between  $s \in \InitSolSet_{\ell}$ and $\MultiGroup_i \in \MultiGroupSet$ if $d(s, \MultiGroup_i) \le \lambda$. Furthermore, for a fixed $\ell \in [k]$, we define $\lambda_{\ell}$ as the minimum $\lambda$ such that there exists a \emph{left-perfect matching} on $H^{\ell}_{\lambda}$. 
\end{definition}

By using the projection graph, we can obtain a $\gamma+1$-distortion feasible (fair) solution from a progressive $\gamma$-cover $T$~\cite{gadekar2025fair,jones2020fair}.



























